% GENERATED FILE. Edit canonical lessons, then run npm run generate:books.
% canonical-source-hash: fnv1a64:413ec22a7d9f31aa
% canonical-lessons: MW-C253-sardi, MW-C253-khasi, MW-C253-jivno, MW-C253-dharakno, MW-C253-dukhno, MW-R253-review-248-253

\chapter{Colds, Coughs and a Long Life}
\label{ch:mw-colds-coughs-and-a-long-life}
\hlchaptermodality{253}

\begin{chapteropening}
\textbf{By the end of this chapter:} \emph{I can say what is wrong when I feel ill.}

Complete the last lesson of chapter 253: I can say what is wrong when I feel ill.
\end{chapteropening}

\section[\texorpdfstring{\mw{सरदी} (sardī)}{सरदी (sardī)}]{\mw{सरदी} (sardī) — a cold (the illness)}

\label{lesson:MW-C253-sardi}



\begin{quote}
Before the new one: say the Marwadi for to chew, then the Marwadi for an ache.
\end{quote}

\subsection*{You'll want to know: \mw{सरदी}}
\textbf{\mw{सरदी}} — \emph{sardī} — \textquotedblleft{}a cold (the illness)\textquotedblright{}.

\begin{quote}
\mw{बेटो} \mw{बीमार} \mw{है}, \mw{सरदी} \mw{है।}
\end{quote}

\begin{quote}
\emph{beṭo bīmār hai, sardī hai.}
\end{quote}

\begin{quote}
My son is ill; he has a cold.
\end{quote}

A feminine noun.

\begin{cousinweb}
\textbf{\mw{सरदी}} (\emph{sardī}) is usually traced to Persian \emph{sardī}, \textquotedblleft{}coldness\textquotedblright{}; in Marwari it names the illness.
\end{cousinweb}

\begin{grammarlens}[title={What you've built}]
One more word for this chapter.
\end{grammarlens}

\subsection*{Your turn}
\begin{itemize}
\raggedright
  \item \emph{Say it:} \emph{sardī}
  \item \emph{Say it:} \emph{sardī}, once more
  \item \emph{Say it:} \emph{beṭo bīmār hai, sardī hai.} — My son is ill; he has a cold.
  \item \emph{Recall:} say the Marwadi for to chew, then the Marwadi for an ache, then say \emph{sardī} again
\end{itemize}

\begin{tcolorbox}[breakable,colback=teal!4,colframe=teal!35!black,title={Before you move on}]
What does \mw{सरदी} mean? (\textquotedblleft{}A cold (the illness)\textquotedblright{}.) Say it once more.
\end{tcolorbox}
\section[\texorpdfstring{\mw{खांसी} (khā̃sī)}{खांसी (khā̃sī)}]{\mw{खांसी} (khā̃sī) — a cough}

\label{lesson:MW-C253-khasi}



\begin{quote}
Before the new one: say the Marwadi for an ache, then the Marwadi for a cold (the illness).
\end{quote}

\subsection*{You'll want to know: \mw{खांसी}}
\textbf{\mw{खांसी}} — \emph{khā̃sī} — \textquotedblleft{}a cough\textquotedblright{}.

\begin{quote}
\mw{रात} \mw{मांय} \mw{घणी} \mw{खांसी} \mw{आवै} \mw{है।}
\end{quote}

\begin{quote}
\emph{rāt mā̃y ghaṇī khā̃sī āvai hai.}
\end{quote}

\begin{quote}
I cough a lot at night.
\end{quote}

A feminine noun.

\begin{cousinweb}
\textbf{\mw{खांसी}} (\emph{khā̃sī}) is inherited from Sanskrit \emph{kāsā}, \textquotedblleft{}a cough\textquotedblright{}. CDIAL gives no Marwari form in this entry; it lists Hindi and Gujarati \emph{khā̃sī} (CDIAL 3138).
\end{cousinweb}

\begin{grammarlens}[title={What you've built}]
One more word for this chapter.
\end{grammarlens}

\subsection*{Your turn}
\begin{itemize}
\raggedright
  \item \emph{Say it:} \emph{khā̃sī}
  \item \emph{Say it:} \emph{khā̃sī}, once more
  \item \emph{Say it:} \emph{rāt mā̃y ghaṇī khā̃sī āvai hai.} — I cough a lot at night.
  \item \emph{Recall:} say the Marwadi for an ache, then the Marwadi for a cold (the illness), then say \emph{khā̃sī} again
\end{itemize}

\begin{tcolorbox}[breakable,colback=teal!4,colframe=teal!35!black,title={Before you move on}]
What does \mw{खांसी} mean? (\textquotedblleft{}A cough\textquotedblright{}.) Say it once more.
\end{tcolorbox}
\section[\texorpdfstring{\mw{जीवणो} (jīvṇo)}{जीवणो (jīvṇo)}]{\mw{जीवणो} (jīvṇo) — to be alive; to live (a span of years)}

\label{lesson:MW-C253-jivno}



\begin{quote}
Before the new one: say the Marwadi for a cold (the illness), then the Marwadi for a cough.
\end{quote}

\subsection*{You'll want to know: \mw{जीवणो}}
\textbf{\mw{जीवणो}} — \emph{jīvṇo} — \textquotedblleft{}to be alive; to live (a span of years)\textquotedblright{}.

\begin{quote}
\mw{म्हारो} \mw{काको} \mw{अस्सी} \mw{बरस} \mw{जीयो।}
\end{quote}

\begin{quote}
\emph{mhāro kāko assī baras jīyo.}
\end{quote}

\begin{quote}
My father's brother lived to eighty.
\end{quote}

Past (he …): \mw{वो} \mw{सो} \mw{बरस} \mw{जीयो}, \emph{vo so baras jīyo}.

\begin{cousinweb}
\textbf{\mw{जीवणो}} (\emph{jīvṇo}) is inherited from Sanskrit \emph{jīvati}, \textquotedblleft{}is alive\textquotedblright{}; CDIAL lists Old Marwari \emph{jīvaï}, \textquotedblleft{}to live\textquotedblright{}, under it (CDIAL 5241).
\end{cousinweb}

\begin{grammarlens}[title={What you've built}]
One more word for this chapter.
\end{grammarlens}

\subsection*{Your turn}
\begin{itemize}
\raggedright
  \item \emph{Say it:} \emph{jīvṇo}
  \item \emph{Say it:} \emph{jīvṇo}, once more
  \item \emph{Say it:} \emph{mhāro kāko assī baras jīyo.} — My father's brother lived to eighty.
  \item \emph{Recall:} say the Marwadi for a cold (the illness), then the Marwadi for a cough, then say \emph{jīvṇo} again
\end{itemize}

\begin{tcolorbox}[breakable,colback=teal!4,colframe=teal!35!black,title={Before you move on}]
What does \mw{जीवणो} mean? (\textquotedblleft{}To be alive; to live (a span of years)\textquotedblright{}.) Say it once more.
\end{tcolorbox}
\section[\texorpdfstring{\mw{धड़कणो} (dhaṛakṇo)}{धड़कणो (dhaṛakṇo)}]{\mw{धड़कणो} (dhaṛakṇo) — to throb; to pound (of the heart)}

\label{lesson:MW-C253-dharakno}



\begin{quote}
Before the new one: say the Marwadi for a cough, then the Marwadi for to be alive.
\end{quote}

\subsection*{You'll want to know: \mw{धड़कणो}}
\textbf{\mw{धड़कणो}} — \emph{dhaṛakṇo} — \textquotedblleft{}to throb; to pound (of the heart)\textquotedblright{}.

\begin{quote}
\mw{म्हारो} \mw{दिल} \mw{घणो} \mw{धड़कै} \mw{है।}
\end{quote}

\begin{quote}
\emph{mhāro dil ghaṇo dhaṛakai hai.}
\end{quote}

\begin{quote}
My heart is pounding hard.
\end{quote}

Past (it …): \mw{वो} \mw{दिल} \mw{धड़क्यो}, \emph{vo dil dhaṛakyo}.

\begin{cousinweb}
\textbf{\mw{धड़कणो}} (\emph{dhaṛakṇo}) is inherited from a reconstructed \emph{*dhaḍ}, \textquotedblleft{}throb, beat\textquotedblright{}. CDIAL gives no Marwari form in this entry; it lists Hindi \emph{dhaṛaknā}, \textquotedblleft{}to throb\textquotedblright{} (CDIAL 6711).
\end{cousinweb}

\begin{grammarlens}[title={What you've built}]
One more word for this chapter.
\end{grammarlens}

\subsection*{Your turn}
\begin{itemize}
\raggedright
  \item \emph{Say it:} \emph{dhaṛakṇo}
  \item \emph{Say it:} \emph{dhaṛakṇo}, once more
  \item \emph{Say it:} \emph{mhāro dil ghaṇo dhaṛakai hai.} — My heart is pounding hard.
  \item \emph{Recall:} say the Marwadi for a cough, then the Marwadi for to be alive, then say \emph{dhaṛakṇo} again
\end{itemize}

\begin{tcolorbox}[breakable,colback=teal!4,colframe=teal!35!black,title={Before you move on}]
What does \mw{धड़कणो} mean? (\textquotedblleft{}To throb; to pound (of the heart)\textquotedblright{}.) Say it once more.
\end{tcolorbox}
\section[\texorpdfstring{\mw{दूखणो} (dūkhṇo)}{दूखणो (dūkhṇo)}]{\mw{दूखणो} (dūkhṇo) — to ache; to hurt}

\label{lesson:MW-C253-dukhno}



\begin{quote}
Before the new one: say the Marwadi for to be alive, then the Marwadi for to throb.
\end{quote}

\subsection*{You'll want to know: \mw{दूखणो}}
\textbf{\mw{दूखणो}} — \emph{dūkhṇo} — \textquotedblleft{}to ache; to hurt\textquotedblright{}.

\begin{quote}
\mw{म्हारो} \mw{माथो} \mw{दूखै} \mw{है।}
\end{quote}

\begin{quote}
\emph{mhāro mātho dūkhai hai.}
\end{quote}

\begin{quote}
My head aches.
\end{quote}

Past (it …): \mw{वो} \mw{दांत} \mw{दूख्यो}, \emph{vo dānt dūkhyo}.

\begin{cousinweb}
\textbf{\mw{दूखणो}} (\emph{dūkhṇo}) is inherited from Sanskrit \emph{duḥkhati}, \textquotedblleft{}pains\textquotedblright{}. CDIAL gives no Marwari form in this entry; it lists Hindi \emph{dūkhnā} (CDIAL 6376).
\end{cousinweb}

\begin{grammarlens}[title={What you've built}]
One more word for this chapter.
\end{grammarlens}

\subsection*{Your turn}
\begin{itemize}
\raggedright
  \item \emph{Say it:} \emph{dūkhṇo}
  \item \emph{Say it:} \emph{dūkhṇo}, once more
  \item \emph{Say it:} \emph{mhāro mātho dūkhai hai.} — My head aches.
  \item \emph{Recall:} say the Marwadi for to be alive, then the Marwadi for to throb, then say \emph{dūkhṇo} again
\end{itemize}

\begin{tcolorbox}[breakable,colback=teal!4,colframe=teal!35!black,title={Before you move on}]
What does \mw{दूखणो} mean? (\textquotedblleft{}To ache; to hurt\textquotedblright{}.) Say it once more.
\end{tcolorbox}
\section[(dialogue)]{Looking back: from Henna and Wedding Finery to Colds, Coughs and a Long Life}

\label{lesson:MW-R253-review-248-253}



\begin{quote}
Say the Marwadi for henna and to marry off.
\end{quote}

\subsection*{Your turn}
Say these aloud:

\begin{itemize}
\raggedright
  \item henna, to marry off, a dowry, to dress up, a bridegroom's garland
  \item Holi, coloured powder for Holi, a good omen, Diwali, worship
  \item to celebrate, festive dishes, festive joy, the waving of a lamp in worship, an uproar
  \item the neck, a shoulder, the chest, an elbow, the waist
  \item a bone, breath, to pant, to chew, an ache
  \item a cold (the illness), a cough, to be alive, to throb, to ache
\end{itemize}

\begin{tcolorbox}[breakable,colback=teal!4,colframe=teal!35!black,title={Before you move on}]
Henna? (\textbf{\mw{मेंहदी}}, \emph{mẽhdī}.) To marry off? (\textbf{\mw{परणावणो}}, \emph{parṇāvṇo}.) A dowry? (\textbf{\mw{दायजो}}, \emph{dāyjo}.) To dress up? (\textbf{\mw{सजणो}}, \emph{sajṇo}.) A bridegroom's garland? (\textbf{\mw{सेहरो}}, \emph{sehro}.) Holi? (\textbf{\mw{होळी}}, \emph{hoḷī}.) Coloured powder for Holi? (\textbf{\mw{गुलाल}}, \emph{gulāl}.) A good omen? (\textbf{\mw{सुगन}}, \emph{sugan}.) Diwali? (\textbf{\mw{दिवाळी}}, \emph{divāḷī}.) Worship? (\textbf{\mw{पूजा}}, \emph{pūjā}.) To celebrate? (\textbf{\mw{मनावणो}}, \emph{manāvṇo}.) Festive dishes? (\textbf{\mw{पकवान}}, \emph{pakvān}.) Festive joy? (\textbf{\mw{उछाव}}, \emph{uchhāv}.) The waving of a lamp in worship? (\textbf{\mw{आरती}}, \emph{ārtī}.) An uproar? (\textbf{\mw{रोळो}}, \emph{roḷo}.) The neck? (\textbf{\mw{गरदन}}, \emph{gardan}.) A shoulder? (\textbf{\mw{कांधो}}, \emph{kā̃dho}.) The chest? (\textbf{\mw{छाती}}, \emph{chhātī}.) An elbow? (\textbf{\mw{कोहणी}}, \emph{kohṇī}.) The waist? (\textbf{\mw{कमर}}, \emph{kamar}.) A bone? (\textbf{\mw{हाड}}, \emph{hāḍ}.) Breath? (\textbf{\mw{सांस}}, \emph{sā̃s}.) To pant? (\textbf{\mw{हांफणो}}, \emph{hā̃phṇo}.) To chew? (\textbf{\mw{चाबणो}}, \emph{chābṇo}.) An ache? (\textbf{\mw{पीड़}}, \emph{pīṛ}.) A cold (the illness)? (\textbf{\mw{सरदी}}, \emph{sardī}.) A cough? (\textbf{\mw{खांसी}}, \emph{khā̃sī}.) To be alive? (\textbf{\mw{जीवणो}}, \emph{jīvṇo}.) To throb? (\textbf{\mw{धड़कणो}}, \emph{dhaṛakṇo}.) To ache? (\textbf{\mw{दूखणो}}, \emph{dūkhṇo}.)
\end{tcolorbox}
